\section*{अध्याय \ref{ch:g11:fossil-fuel-geology} --- कोयले, तेल और गैस का भूविज्ञान}

\begin{solution}{exo:g11:fossil-fuel-geology:1}
मृत प्लवक ऑक्सीजन-विहीन समुद्री तल पर बैठते हैं; कार्बनिक पदार्थ
कीचड़ में \omterm{def:g11:fossil-fuel-geology:source-rock}{केरोजन} बन जाता है; नए अवसादों के नीचे दफ़न; \omterm{def:g11:fossil-fuel-geology:oil-window}{तेल-खिड़की} में
\omterm{def:g11:fossil-fuel-geology:oil-window}{परिपक्वन}; \omterm{def:g11:fossil-fuel-geology:migration}{प्रवासन}; तेल \omterm{def:g11:fossil-fuel-geology:reservoir}{आवरण चट्टान} के नीचे इकट्ठा होता है।
\end{solution}

\begin{solution}{exo:g11:fossil-fuel-geology:2}
(1) \omterm{def:g11:fossil-fuel-geology:source-rock}{स्रोत चट्टान}। (2) \omterm{def:g11:fossil-fuel-geology:reservoir}{भंडार चट्टान}। (3) \omterm{def:g11:fossil-fuel-geology:reservoir}{आवरण चट्टान}। (4) \omterm{def:g11:fossil-fuel-geology:reservoir}{आवरण चट्टान}। (5) \omterm{def:g11:fossil-fuel-geology:reservoir}{भंडार चट्टान}।
\end{solution}

\begin{solution}{exo:g11:fossil-fuel-geology:3}
शीर्ष: $(60 - 10)/25 = \qty{2.0}{km}$; आधार: $(120 - 10)/25 = \qty{4.4}{km}$।
खिड़की $60/25 = \qty{2.4}{km}$ मोटी है।
\end{solution}

\begin{solution}{exo:g11:fossil-fuel-geology:4}
(1) पीट, लिग्नाइट, बिटुमिनी \omterm{def:g11:fossil-fuel-geology:coal}{कोयला}, ऐन्थ्रेसाइट। (2) ऐन्थ्रेसाइट, 86 से
\qty{97}{\percent} कार्बन। (3) ऐन्थ्रेसाइट, लगभग \qty{5}{km} और
\qty{150}{\degreeCelsius} से आगे।
\end{solution}

\begin{solution}{exo:g11:fossil-fuel-geology:5}
(1) तरल घनत्व के अनुसार स्वयं को छाँट लेते हैं: गैस सबसे हल्की है, फिर तेल,
फिर पानी। (2) तेल और गैस इसलिए ऊपर उठते हैं कि वे पानी से हल्के हैं; केवल
वह चट्टान उन्हें रोकती है जिससे होकर वे नहीं गुज़र सकते। पंकाश्म (या शेल) और लवण।
\end{solution}

\begin{solution}{exo:g11:fossil-fuel-geology:6}
(1) A: $(60 - 15)/20 = \qty{2.25}{km}$ और $(120 - 15)/20 = \qty{5.25}{km}$।
B: $45/40 \approx \qty{1.13}{km}$ और $105/40 \approx \qty{2.63}{km}$।
(2) \qty{2}{km} पर: A में $15 + 40 = \qty{55}{\degreeCelsius}$, अभी
परिपक्व नहीं, कोई तेल नहीं; B में $15 + 80 = \qty{95}{\degreeCelsius}$, \omterm{def:g11:fossil-fuel-geology:oil-window}{तेल-खिड़की}
में, वह तेल बनाती है। (3) द्रोणी A में: वहाँ तेल केवल
\qty{2.25}{km} से नीचे बनता है।
\end{solution}

\begin{solution}{exo:g11:fossil-fuel-geology:7}
(1) $11 \times 3 = \qty{33}{m}$ पीट। (2) प्रति वर्ष \qty{1}{mm} पर
$\qty{33000}{mm}$: लगभग \num{33000} वर्ष। (3) पीट केवल वहीं बचा रहता है जहाँ
वह ठहरे पानी के नीचे, ऑक्सीजन की पहुँच से बाहर हो; जैसे-जैसे वह ढेर होता है,
भूमि को उसी गति से धँसना चाहिए ताकि उसकी सतह जलमग्न रहे।
\end{solution}

\begin{solution}{exo:g11:fossil-fuel-geology:8}
(1) $20 + 30 \times 1.2 = \qty{56}{\degreeCelsius}$;
$20 + 30 \times 2.3 = \qty{89}{\degreeCelsius}$;
$20 + 30 \times 4.0 = \qty{140}{\degreeCelsius}$।
(2) 120 मिलियन वर्ष पहले अपरिपक्व; 95 मिलियन वर्ष पहले \omterm{def:g11:fossil-fuel-geology:oil-window}{तेल-खिड़की}
में; 20 मिलियन वर्ष पहले \omterm{def:g11:fossil-fuel-geology:oil-window}{गैस-खिड़की} में।
\end{solution}

\begin{solution}{exo:g11:fossil-fuel-geology:9}
(1) $2580 - 2400 = \qty{180}{m}$। (2) गैस
$2450 - 2400 = \qty{50}{m}$ भरती है, तेल $2580 - 2450 = \qty{130}{m}$। (3) वे सबसे गहरी बंद
समोच्च रेखा के नीचे, अतिप्रवाह बिंदु पर, बह निकलते हैं, और ऊपर की ओर आगे \omterm{def:g11:fossil-fuel-geology:migration}{प्रवासन} करते हैं,
किसी दूसरे \omterm{def:g11:fossil-fuel-geology:trap}{ट्रैप} तक या सतह तक।
\end{solution}

\begin{solution}{exo:g11:fossil-fuel-geology:10}
(1) \qty{3}{km} पर खिड़की का शीर्ष: $(60 - 20)/3 \approx
\qty{13}{\degreeCelsius/km}$; आधार: $(120 - 20)/3 \approx
\qty{33}{\degreeCelsius/km}$। (2) \qty{3}{km} पर तापमान
$20 + 75 = \qty{95}{\degreeCelsius}$ से $20 + 102 = \qty{122}{\degreeCelsius}$ तक है:
\omterm{def:g11:fossil-fuel-geology:oil-window}{तेल-खिड़की} में, या सबसे तीव्र प्रवणताओं के लिए उसके आधार पर। नियम
उचित है: अधिकांश द्रोणियों में \qty{3}{km} चट्टानों के नीचे की \omterm{def:g11:fossil-fuel-geology:source-rock}{स्रोत चट्टान}
तेल बनाना शुरू कर चुकी होती है।
\end{solution}

\begin{solution}{exo:g11:fossil-fuel-geology:11}
(1) \qty{30}{\percent}: लिग्नाइट; \qty{70}{\percent}: बिटुमिनी \omterm{def:g11:fossil-fuel-geology:coal}{कोयला};
\qty{92}{\percent}: ऐन्थ्रेसाइट। (2) ऐन्थ्रेसाइट, लगभग
\qty{5}{km} और \qty{150}{\degreeCelsius} से आगे। (3) काला और काँच जैसा।
\end{solution}

\begin{solution}{exo:g11:fossil-fuel-geology:12}
(1) $1100/25 = \qty{44}{m}$ प्रति मिलियन वर्ष। (2) $0.044 \times 30
\approx \qty{1.3}{\degreeCelsius}$ प्रति मिलियन वर्ष। (3) खिड़की
\qty{60}{\degreeCelsius} में फैली है: $60/1.32 \approx 45$ मिलियन वर्ष। (4) 95 मिलियन वर्ष
पहले के बाद दफ़न धीमा हो गया (लगभग 24, फिर 23 और
\qty{20}{m} प्रति मिलियन वर्ष): \omterm{def:g11:fossil-fuel-geology:source-rock}{स्रोत चट्टान} अधिक धीरे गरम हुई।
\end{solution}

\begin{solution}{exo:g11:fossil-fuel-geology:13}
(1) $(80 - 51)/(117 - 51) = 29/66 \approx 0.44$: लगभग \qty{44}{\percent}।
(2) ऊपर किसी \omterm{def:g11:fossil-fuel-geology:trap}{ट्रैप} के होने से पहले ही वह ऊपर की ओर प्रवासित हो गया: वह
सतह पर निकल गया, जहाँ ऑक्सीकृत हुआ, या किसी दूसरे \omterm{def:g11:fossil-fuel-geology:trap}{ट्रैप} में इकट्ठा हुआ। (3) 40
मिलियन वर्ष पहले तक \omterm{def:g11:fossil-fuel-geology:source-rock}{स्रोत चट्टान} \omterm{def:g11:fossil-fuel-geology:oil-window}{तेल-खिड़की} छोड़ चुकी थी: वह केवल गैस बनाती थी,
और अपनति उससे केवल गैस रख सकती थी।
\end{solution}

\begin{solution}{exo:g11:fossil-fuel-geology:14}
(1) शीर्ष: $(60 - 20)/40 = \qty{1.0}{km}$; आधार: $(120 - 20)/40 =
\qty{2.5}{km}$। (2) \qty{1.0}{km} \qty{0.4}{km} (150 मिलियन
वर्ष पहले) और \qty{1.2}{km} (120) के बीच है: $150 - 30 \times 0.6/0.8 = 127.5$
मिलियन वर्ष पहले। \qty{2.5}{km} \qty{2.3}{km} (95) और
\qty{3.0}{km} (66) के बीच है: $95 - 29 \times 0.2/0.7 \approx 86.7$ मिलियन वर्ष
पहले। (3) कोई तेल नहीं: \omterm{def:g11:fossil-fuel-geology:source-rock}{स्रोत चट्टान} अपनति बनने से पहले ही \omterm{def:g11:fossil-fuel-geology:oil-window}{तेल-खिड़की} छोड़ चुकी थी,
और उसका तेल प्रवासित होकर जा चुका था। वह 80 मिलियन
वर्ष पहले के बाद बनी गैस रख सकती थी।
\end{solution}

\begin{solution}{exo:g11:fossil-fuel-geology:15}
(1) कम से कम \qty{5}{km}। (2) $15 + 30 \times 5 = \qty{165}{\degreeCelsius}$,
\qty{150}{\degreeCelsius} से ऊपर: संगत। (3) कम से कम \qty{5}{km}
चट्टान उठाकर अपरदित की गई है। अधिक से अधिक \num{299} मिलियन वर्षों में
$\qty{5e6}{mm}$ औसतन कम से कम $5 \times 10^6/2.99 \times 10^8
\approx \qty{0.017}{mm}$ प्रति वर्ष देता है।
\end{solution}

\begin{solution}{pb:g11:fossil-fuel-geology:1}
\textbf{1.} \omterm{def:g11:fossil-fuel-geology:source-rock}{स्रोत चट्टान}: काला शेल; भंडार: सरंध्र बलुआ पत्थर;
\omterm{def:g11:fossil-fuel-geology:reservoir}{आवरण चट्टान}: उसके ऊपर का पंकाश्म।

\textbf{2.} ऑक्सीजन-विहीन समुद्री तल पर बैठे प्लवक सड़ने से बच गए
और कीचड़ में दब गए।

\textbf{3.} \omterm{def:g11:fossil-fuel-geology:source-rock}{स्रोत चट्टान} जुरैसिक में (165, 201.4 और
143.1 के बीच है)। बलुआ पत्थर जुरैसिक में शुरू हुआ और ज़्यादातर
पूर्व क्रिटेशस में जमा।

\textbf{4.} किसी अवसादी चट्टान में बिखरा कार्बनिक पदार्थ, जो
कार्बनिक विलायकों में अघुलनशील है। वह लगभग
\qty{60}{\degreeCelsius} तक गरम होने पर ही तेल बनाता है, जिसके लिए दफ़न चाहिए।

\textbf{5.} $400 + 800 + 1100 + 700 + 700 + 300 - 500 = \qty{3500}{m}$।

\textbf{6.} \qty{0.4}{km}, \qty{1.2}{km}, \qty{2.3}{km}, \qty{3.0}{km},
\qty{3.7}{km}, \qty{4.0}{km}; आज \qty{3.5}{km}।

\textbf{7.} $400/15 \approx 26.7$; $800/30 \approx 26.7$; $1100/25 = 44$;
$700/29 \approx 24.1$; $700/31 \approx 22.6$; $300/15 = 20$ (\unit{m} प्रति
मिलियन वर्ष)।

\textbf{8.} 20 मिलियन वर्ष पहले, \qty{4.0}{km} पर।

\textbf{9.} द्रोणी ऊपर उठी और उसके शीर्ष से \qty{500}{m} चट्टान
अपरदित हुई।

\textbf{10.} बलुआ पत्थर का आधार 150 मिलियन वर्ष पहले, शेल और पंकाश्म के
\qty{400}{m} के ऊपर, जमा था: $3.5 - 0.4 = \qty{3.1}{km}$।

\textbf{11.} $T = 20 + 30 z$: \qty{32}{\degreeCelsius}, \num{56}, \num{89},
\num{110}, \num{131}, \qty{140}{\degreeCelsius}; आज
\qty{125}{\degreeCelsius}।

\textbf{12.} $40/30 \approx \qty{1.33}{km}$ और $100/30 \approx
\qty{3.33}{km}$।

\textbf{13.} 120 मिलियन वर्ष पहले (\qty{56}{\degreeCelsius}) और
95 मिलियन वर्ष पहले (\qty{89}{\degreeCelsius}) के बीच।

\textbf{14.} तापमान 25 मिलियन वर्षों में 33 डिग्री बढ़ा;
\num{56} से \qty{60}{\degreeCelsius} तक 4 डिग्री बढ़ने में उसे $25 \times 4/33 \approx 3.0$ मिलियन वर्ष लगे: लगभग $120 - 3.0 = 117$ मिलियन वर्ष
पहले।

\textbf{15.} 66 मिलियन वर्ष पहले (\qty{110}{\degreeCelsius}) और
35 (\qty{131}{\degreeCelsius}) के बीच: $66 - 31 \times 10/21 \approx 51.2$ मिलियन
वर्ष पहले।

\textbf{16.} \qty{140}{\degreeCelsius}, 20 मिलियन वर्ष पहले। आज
\qty{125}{\degreeCelsius} पर वह \omterm{def:g11:fossil-fuel-geology:oil-window}{गैस-खिड़की} में है: वह ज़्यादातर
मीथेन बनाती है।

\textbf{17.} अपने ऊपर कोई \omterm{def:g11:fossil-fuel-geology:trap}{ट्रैप} न होने से वह ऊपर प्रवासित होकर निकल गया,
सतह तक या दूसरे \omterm{def:g11:fossil-fuel-geology:trap}{ट्रैपों} तक।

\textbf{18.} $(80 - 51.2)/(117.0 - 51.2) \approx 0.44$।

\textbf{19.} सबसे ऊपर गैस, फिर तेल, फिर पानी।

\textbf{20.} \omterm{def:g11:fossil-fuel-geology:source-rock}{स्रोत चट्टान} ने लगभग \textbf{117
मिलियन वर्ष पहले} \omterm{def:g11:fossil-fuel-geology:oil-window}{तेल-खिड़की} में प्रवेश किया।
\end{solution}
